\documentclass[12pt]{article}
\usepackage[utf8]{inputenc}
\usepackage[T1]{fontenc}
\usepackage{hyperref}
\title{Umiida Language Resources for K-12 Curriculum Development}
\author{Compiled by Claude Sonnet 5.5 (claude-sonnet-5-5) for Andr\'as Kornai}
\date{}
\begin{document}
\maketitle

\section{Introduction}
Umiida (also Umida, Umede; ISO 639-3 xud) is a dormant Aboriginal Australian language of the Kimberley region of Western Australia, with territory along the Yampi Sound coast. Sources class it as Western Worrorran (non-Pama-Nyungan). No remaining community members use it as a first language. My searches found no school materials of any kind. Everything listed below is reference material about the language, not curriculum material.

\section{Primary Educational Resources}
I found no primers, textbooks, or school programmes for Umiida. A K-12 curriculum would have to be built from scratch on the sparse documentation listed below.

\section{Linguistic and Reference Materials}
\begin{itemize}
\item AIATSIS AUSTLANG entry K49 for Umiida. It places the language in Western Worrorran and notes that McGregor and Coate collected vocabulary and sentences. \url{https://collection.aiatsis.gov.au/austlang/language/K49}
\item Wikipedia article on Umiida. It says that what little is known of the language was recorded by Howard Coate in the 1960s. \url{https://en.wikipedia.org/wiki/Umiida}
\item Grambank entry for Umiida (umii1235). \url{https://grambank.clld.org/languages/umii1235}
\item University of Sydney archives record that appeared in the search. I did not open it, so its contents are unverified. \url{https://archives-search.sydney.edu.au/nodes/view/18773}
\end{itemize}

\section{Religious and Community Resources}
I found no scripture translations and no community organisation specific to Umiida.

\section{Digital and Audio Resources}
\begin{itemize}
\item Endangered Languages Project catalogue entry. It lists Umiida among languages that may have no current speakers but whose descendants may still identify with them. \url{https://endangeredlanguages.com/lang/6725?hl=es}
\end{itemize}
I found no audio or video recordings in the search results.

\section{Curriculum Development Context}
Umiida is dormant and no school teaches it. Any curriculum would be a revival or awareness curriculum. A sensible first step is to contact AIATSIS and the Kimberley Aboriginal language centres about archived Coate and McGregor materials. Descendant communities should decide how, and whether, the material is used.

\end{document}